\mainlanguage[de]

\setuplayout
  [backspace=2.0cm,
   height=middle,
   width=middle,
   leftmargin=1.8cm,
   footerdistance=0.0cm,
   rightmargindistance=0.2cm,
   rightmargin=1.8cm,
   bottom=0.5cm,
   leftmargindistance=0.2cm,
   header=2.3cm,
   headerdistance=0.5cm,
   footer=0.5cm,
   top=0cm,
   topspace=0.5cm,
   bottomspace=0.0cm]

\setupformulas[align=right]

\definemeasure[headerimageheight][\dimexpr(\headerheight-2ex)]

\definelayer[sabel]

\startsetups[sabel]

\setlayerframed[sabel][x=0mm,y=0mm][align=flushright, width=\dimexpression{\paperwidth-3cm}, frame=off]{\externalfigure[../include/MAB.jpg][width=3cm]}

\stopsetups

\setupbackgrounds[header][text][bottomframe=off,setups=sabel,background=sabel]
\setuppagenumbering[location={footer,middle},style=\bfc]

\setupfootertexts
  [BE: 21 Arbeitszeit: 30min][Viel Erfolg!]
  [Viel Erfolg!][BE: 21 Arbeitszeit: 30min]

\setdigitmode 3

\setupitemgroup[itemize][1][n,repeat][width=1em]
\setupitemgroup[itemize][2][n,repeat][width=2em]
\setupitemgroup[itemize][3][n,repeat][width=3em]
\setupitemgroup[itemize][4][n,repeat][width=4em]

\definefont[calibribold][name:calibrib at 14pt]
\definefont[calibri][name:calibri at 11pt]

\starttext
%\showlayout
%\setupheader[state=stop]
\setupinterlinespace[line=4.2ex]
% --- School name (Calibri, 14pt, bold, centered)
{
\calibribold
\centerline{Fachakademie für Wirtschaft der Mittelstandsakademie Bayern}
\centerline{1.Kurzarbeit im Wahlpflichtfach Mathematik FHR}
}

\blank[big]

{
\calibri
\starttabulate[|p(9cm)|p(9cm)|]
\NC Aufgabenstellung: Bernd Donner \NC Bearbeitungszeit: 30min \NC\NR
\NC Gesamtpunktzahl: 21BE          \NC Hilfsmittel: TR, MH     \NC\NR
\stoptabulate

\blank[big]

\starttabulate[|p(6cm)|p(6cm)|p(6cm)|]
\NC Name: {\vrule width 5cm height 0mm depth .5pt}
\NC Datum: 11.11.2025
\NC Klasse: VZ2 \NC\NR

\stoptabulate
}
\page[yes]
%\setupheader[state=start]
\setupinterlinespace[line=2.8ex]

Hinweis: Achten Sie auf eine übersichtliche und schöne äußere Form,
korrekte mathematische Schreibweise und geben Sie bei Textaufgaben kurze
Antworten an.

\startitemize
\item Bestimmen Sie die Ableitung für folgende Funktionen:

\startitemize[columns, two, packed][grid=no, inner={\setupinterlinespace[24pt]}]
      \startitem $f(x) = \frac{2x+1}{2x-1}$  (BE~2) \stopitem
      \startitem $g(x) = \sin(\frac{1}{x})$  (BE~2) \stopitem
      \startitem $h(x) = \sqrt{x}$           (BE~2) \stopitem
\stopitemize

\item In einer großen Stadt breitet sich eine Viruserkrankung aus.
Die momentane Erkrankungsrate wird modellhaft beschrieben durch die
Funktion \math{f} mit
\startformula
f(t) = 150 t^2 e^{-0.2 t} \qquad\text{für}\qquad t \geq 0
\stopformula
Dabei ist \math{t} die Zeit in Wochen seit Beobachtungsbeginn und \math{f(t)}
die Anzahl der Neuerkrankungen pro Woche.

\startitemize
\item Zeichnen Sie den Verlauf der Neuerkrankungen zwischen dem
Beobachtungsbeginn und der 60sten Woche. (BE~6)
\item In welcher Woche ist nach diesem Modell mit der größten Zahl an
Neuerkrankungen zu rechnen. Beantworten Sie diese Frage wie im Unterricht
behandelt oder durch einen gleichwertigen exakten mathematischen Nachweis. (BE~9)
\stopitemize
\stopitemize
\stoptext

